\section*{表目录（原书页码）}\phantomsection\label{front:tables}
\addcontentsline{toc}{section}{表目录（原书页码）}
\begingroup\small
\begin{longtable}{@{}p{.86\textwidth}r@{}}
\hyperref[tab:1.2.1]{1.2.1 排列方式的数目} & \hyperref[tab:1.2.1]{13} \\
\hyperref[tab:4.1.1]{4.1.1 联合概率质量函数$f(x,y)$的取值} & \hyperref[tab:4.1.1]{115} \\
\hyperref[tab:7.3.1]{7.3.1 二项分布参数$p$的三个估计量} & \hyperref[tab:7.3.1]{286} \\
\hyperref[tab:7.4.2]{7.4.2 白血病病例计数} & \hyperref[tab:7.4.2]{291} \\
\hyperref[tab:8.3.1]{8.3.1 假设检验中的两类错误} & \hyperref[tab:8.3.1]{310} \\
\hyperref[tab:9.2.1]{9.2.1 位置–尺度枢轴量} & \hyperref[tab:9.2.1]{346} \\
\hyperref[tab:9.2.2]{9.2.2 Sterne的接受域与置信集} & \hyperref[tab:9.2.2]{349} \\
\hyperref[tab:9.3.3]{9.3.3 三个90\%正态置信区间} & \hyperref[tab:9.3.3]{357} \\
\hyperref[tab:10.1.1]{10.1.1 自助法与Delta方法的方差} & \hyperref[tab:10.1.1]{387} \\
\hyperref[tab:10.2.2]{10.2.2 中位数相对于均值的渐近相对效率} & \hyperref[tab:10.2.2]{391} \\
\hyperref[tab:10.2.3]{10.2.3 Huber估计量} & \hyperref[tab:10.2.3]{391} \\
\hyperref[tab:10.2.4]{10.2.4 Huber估计量的渐近相对效率，$k=1.5$} & \hyperref[tab:10.2.4]{393} \\
\hyperref[tab:10.3.5]{10.3.5 Poisson似然比检验统计量} & \hyperref[tab:10.3.5]{396} \\
\hyperref[tab:10.3.6]{10.3.6 稳健检验的功效} & \hyperref[tab:10.3.6]{401} \\
\hyperref[tab:10.4.7]{10.4.7 枢轴区间的置信系数} & \hyperref[tab:10.4.7]{404} \\
\hyperref[tab:10.4.8]{10.4.8 基于Huber的M估计量的区间的置信系数} & \hyperref[tab:10.4.8]{407} \\
\hyperref[tab:11.3.1]{11.3.1 图11.3.1中绘制的数据} & \hyperref[tab:11.3.1]{439} \\
\hyperref[tab:11.3.2]{11.3.2 简单线性回归的方差分析表} & \hyperref[tab:11.3.2]{449} \\
\hyperref[tab:12.3.1]{12.3.1 挑战者号数据} & \hyperref[tab:12.3.1]{481} \\
\hyperref[tab:12.4.2]{12.4.2 长鼻袋鼠数据} & \hyperref[tab:12.4.2]{485} \\
\hyperref[tab:12.4.3]{12.4.3 回归M估计量的渐近相对效率} & \hyperref[tab:12.4.3]{487} \\
\hyperref[tab:12.5.4]{12.5.4 按年龄分组的雀鹰存活情况} & \hyperref[tab:12.5.4]{491} \\
\end{longtable}\endgroup
